\section*{1. Ausgangslage / Stand der Forschung}
Die Legacy-Code-Base und der dazugehörige Testpool des Entwicklungsteams sind über die Jahrzehnte gewachsen, veraltet und durch eine starke Verteilung der Umgebung unübersichtlich. Die Ergebnisse des Testpools müssen von verschiedenen Stellen zusammengetragen werden, wodurch das tägliche Monitoring viel Zeit beansprucht.
Um bestehende Produkte zukunftsfähig zu machen, sind Modernisierungen stets gefragt. In dieser Umgebung muss dafür die Übersichtlichkeit gesteigert und der Arbeitsaufwand verringert werden.

Moderne Standards sehen eine übersichtliche, kompakte und nutzerzentrierte Lösung vor. Dabei müssen wichtige Kenngrößen präsent dargestellt sein und auf einem Blick erfasst werden können. Ein konsistentes Design, welches u.a. verwandte Metriken gruppiert, ist ebenso sinnvoll.

Ein neuer Standard trifft auf viele verschiedene Meinungen. Da das Team mit seinem Produkt mitgewachsen ist, kamen bereits vor Projektstart der erste Unmut und Kritik auf. Es wurde die Sorge geäußert, ob die bisherige Lösung nicht mehr ausreicht und diese, trotz Funktionstüchtigkeit, vollständig ersetzt werden soll. Um die Akzeptanz zu steigern und den Übergang zu erleichtern, ist eine nutzernahe Entwicklung entscheidend.

Hierfür wird mithilfe einer kurzen Feedbackschleife gearbeitet, welche in Form eines Interviews umgesetzt wird. Es gibt verschiedene Standards, dieses aufzustellen, wie beispielsweise In-Depth-Interviews (Tiefen- / Intensivinterview) und Semi-strukturierte Interviews. Hier wird ein erprobter Standard ausgewählt und das Interview danach aufgestellt.

% ************************************************************************************
\section*{2. Erkenntnisinteresse / Motivation}
Das Interesse an dem Thema entstand durch die Vorlesungen, bei denen insbesondere der positive Einfluss auf Projektverläufe durch die Anwendung von strukturiertem und wartbarem Code herausstach. Diese theoretischen Erkenntnisse und frischen Ideen sollen mit dem Arbeitsalltag verknüpft werden und einen Beitrag zur Optimierung bestehender Prozesse leisten.

Ziel ist es, herauszufinden, welchen nachhaltigen Mehrwert die vorgeschlagene Dashboard-Lösung im Testpoolmonitoring-Prozess bietet. Um dieses Ziel zu erreichen, werden ebenso Aspekte aus Sicht des Teams, wie Akzeptanz und Nutzerfreundlichkeit, in die Bewertung mit einbezogen.

% ************************************************************************************
\section*{3. Problemstellung / Forschungsfrage}
Die Hauptfragestellung lautet: Wie gut kann ein Dashboard zur Effizienzsteigerung des Testpoolmonitorings beitragen?\\
Dabei sollen folgende Teilfragen beantwortet werden:
\begin{itemize}
    \item Welche Vor- und Nachteile hat die bisherige Lösung?
    \item Welche Herausforderungen bestehen aktuell?
    \item Worin unterscheidet sich die neue Lösung?
    \item Wie lässt sich die Akzeptanz im Team messen?
\end{itemize}

% ************************************************************************************
\section*{4. Ziel der Arbeit}
Als Ziel der Arbeit wird die Steigerung der Effizienz des Testpoolmonitorings unter Einführung eines Dashboards im Vergleich zur bisherigen Lösung analysiert. Dafür werden die Testergebnisse auf diesem Dashboard zusammengefasst dargestellt. Das Produkt wird im Anschluss bewertet, indem die technischen Anforderungen evaluiert und mittels eines Interviews validiert werden.

% ************************************************************************************
\section*{5. Eingrenzung des Themas}
Der Fokus der Arbeit liegt auf dem Monitoring-Prozess. Das Dashboard wird erstellt, um die zu überwachenden Metriken übersichtlich darzustellen. Die Projektskizze beinhaltet dabei den Aufbau und Funktionsweise des Dashboards inkl. des Datenflusses bis zur Anzeige.

Es wird analysiert, inwieweit die neue Lösung die Anforderungen erfüllt und die Effizienz beeinflusst. Weiterhin wird es mithilfe eines Interviews validiert. Hierbei wird die Methodik, bis zur praktischen Umsetzung des Interviews und dessen Auswertung behandelt. Die befragte Personengruppe ist auf das Entwicklungsteam beschränkt, welches das Dashboard nutzen wird.

In der Arbeit wird sich ausschließlich auf diesen Anwendungsfall beschränkt. Das heißt, dass weitere Tools und Strategien zur Qualitätssteigerung und Prozessoptimierung nicht weiter untersucht werden. Der Anbieter zur Dashboard-Erstellung wird nach vorheriger Evaluation ausgewählt. Die restlichen Anbieter sowie alternative Verfahren zur Anzeigenoptimierung werden in diesem Kontext nicht weiter betrachtet. Da das Dashboard die bisherige Lösung in ihrer Funktion ablösen soll, ist auch die Bewertung der Inhalte, wie Vollständigkeit und Sinnhaftigkeit der Tests, nicht Teil der Arbeit.

% ************************************************************************************
\section*{6. Quellen und Materialien}
Verwendete Quellen:
\begin{itemize}
    \item Fachliteratur, Studien
    \item Testpool, Ergebnisanzeige und Umgebung
    \item Firmeninterne Dokumentation
    \item Technische Dokumentationen (Dashboard-Tool etc.)
\end{itemize}

% ************************************************************************************
\section*{7. Methodik}
Nach einer ausführlichen Literaturrecherche folgt die Analyse der Daten und Prozesse, welche beim Monitoring Anwendung finden. Anschließend werden die Anforderungen der bisherigen Lösung und die Erwartungen an die neue Lösung analysiert und festgehalten, um sie zur Evaluation nutzen zu können.

Das Dashboard muss konzipiert und entwickelt werden. Dafür werden verschiedene Tools bewertet und eine Skizze mit den wichtigsten Inhalten und Vorgehensweisen aufgestellt. Daraufhin können das Dashboard anhand der Skizze aufgebaut und die entsprechenden Daten eingepflegt werden.

Die Validierung erfolgt mittels Interviews mit den Teammitgliedern, welche die Lösung verwenden werden. Zur Erstellung des Fragenkatalogs für die Interviews ist eine weitere Literaturrecherche notwendig. Hierfür müssen mögliche Antworten und Ergebnisse vorher in Betracht gezogen werden, um die Fragen auf gewünschte Metriken zu reduzieren. Nach Erstellung des Fragenkatalogs und Fertigstellung des Dashboards werden die Interviews durchgeführt und die Ergebnisse ausgewertet.

% ************************************************************************************
\section*{8. Benötigte Ressourcen}
\begin{itemize}
    \item Entwicklungsumgebung (Laptop mit Docker-fähigem Betriebssystem, IDE, Zugang zu Versionsverwaltungssystem)
    \item Zugang zu Testsystemen und Schnittstellen
    \item Ein Team, welches das die neue Lösung nutzen wird und die Kapazität hat, Interviews und Feedback zu geben
    \item Literatur, Studien, Whitepapers

\end{itemize}
% ************************************************************************************
\section*{9. Vorläufige Gliederung}
\begin{enumerate}
    \item Einleitung
    \item Stand der Forschung
    \item Problemstellung und Zielsetzung
    \item Methodik
    \begin{enumerate}
        \item Evaluation des Dashboard-Tools
        \item Konzeption und Erstellung des Dashboards
        \item Aufbau und Erstellung des Fragenkatalogs für das Interview
    \end{enumerate}
    \item Entwicklung des Grafana Dashboards
    \begin{enumerate}
        \item Evaluation verschiedener Tools
        \item Aufbau und Inhalt
    \end{enumerate}
    \item Ergebnisse
    \begin{enumerate}
        \item Evaluation der technischen Anforderungen
        \item Auswertung des Interviews
    \end{enumerate}
    \item Diskussion
    \item Fazit und Ausblick
\end{enumerate}
% ************************************************************************************
%\newpage
\section*{10. Zeitplan}
\begin{table}[h]
\begin{tabular}{|l|c|l|l|}
\hline
Monat                 & Wo & Bearbeitete Themen & Abgaben / Termine \\ \hline
% -------------------------------------------------------------------------------------
\multirow{3}{*}{Juni} & 1  & Literaturrecherche, Methodikkapitel  &                   \\ \cline{2-4} 
                      & 2  & \makecell[l]{Grundlagen - Stand der Forschung,\\ Auswertung Ergebnisse} &                   \\ \cline{2-4} 
                      & 3  & Aufarbeitung Interviews, Einleitung & 1. Treffen / Absprache \\ \hline
% --------------------------------------------------------------------------------
\multirow{5}{*}{Juli} & 4  & Ergebniskapitel – Aufbau Dashboard &                   \\ \cline{2-4} 
                      & 5  & Ergebniskapitel – Evaluation &                   \\ \cline{2-4} 
                      & 6  & Diskussion - Validierung & \makecell[l]{2. Treffen / Absprache, \\ Abgabe Leseprobe} \\ \cline{2-4} 
                      & 7  & Einarbeitung von Feedback & Besprechung Leseprobe \\ \cline{2-4} 
                      & 8  & Grundlagen - Theorie &                   \\ \hline
% --------------------------------------------------------------------------------
\multirow{4}{*}{Aug}  & 9  & Grundlagen         &                   \\ \cline{2-4} 
                      & 10 & Einleitung / Fazit &                   \\ \cline{2-4} 
                      & 11 &                    &                   \\ \cline{2-4} 
                      & 12 &                    &                   \\ \hline
% --------------------------------------------------------------------------------
\multirow{2}{*}{Sep}  & 13 & Formatierung, Korrektur &                   \\ \cline{2-4} 
                      & 14 & Formatierung, Korrektur, Abgabe & Abgabe \\ \hline
\end{tabular}
\end{table}

% ************************************************************************************
\section*{11. Vorläufige Literatur}
\begin{itemize}
    \item Adeoye‐Olatunde, Omolola A.; Olenik, Nicole L. (2021): Research and scholarly methods: Semi‐structured interviews. In: J Am Coll Clin Pharm 4 (10), S. 1358–1367. DOI: \url{https://doi.org/10.1002/jac5.1441}.
    
    \item Feathers, Michael C. (2013): Working effectively with legacy code. 14. print. Upper Saddle River, N.J.: Prentice Hall (Robert C. Martin series). ISBN: 0131177052.
    
    \item Holtzblatt, Karen (2015): Contextual design. Evolved. San Rafael, California (1537 Fourth Street, San Rafael, CA 94901 USA): Morgan \& Claypool (Synthesis lectures on human-centered informatics, \# 24). DOI: \url{https://doi.org/10.2200/S00597ED1V01Y201409HCI024}.
    
    \item IDEO (2015): The field guide to human-centered design. Design kit. 1st. ed. San Francisco, Calif.: Design Kit. ISBN: 978-0-9914063-1-9.
    
    \item Kaufmann, Traute (2025): Interview. In: Traute Kaufmann (Hg.): Strategiewerkzeuge aus der Praxis Band 1. Analyse und Beurteilung der strategischen Ausgangslage. 2., aktualiserte Auflage. Berlin, Heidelberg: Springer Berlin Heidelberg; Imprint Springer Gabler, S. 69–76. DOI: \url{https://doi.org/10.1007/978-3-662-69887-7_8}.
    
    \item Soares, Marcelo M.; Marcus, Aaron; Rosenzweig, Elizabeth (Hg.) (2021): UX research and design. HCI International. Cham: Springer (Lecture notes in computer science, 12779). DOI: \url{https://doi.org/10.1007/978-3-030-78221-4}.

    \item Williamson, Nicola; Howse, Chloe; Hodson, Nicola; Stein, Julia; Arbuckle, Rob (2025): Qualitative In-trial Interviews: Methods, Challenges, and Best Practice. In: The patient 18 (3), S. 199–209. DOI: \url{https://doi.org/10.1007/s40271-024-00726-2}.
\end{itemize}
